\chapter{Software Documentation}
\label{chap:software}

\section{Code Documentation Standard}

The codebase is implemented in C\# under Unity. All custom scripts adopt the C\#
XML documentation comment standard (triple-slash \texttt{///}) --- the idiomatic
equivalent of JavaDoc for the .NET ecosystem. This standard was chosen because
it is the native format consumed by Visual Studio, Visual Studio Code, and
JetBrains Rider for IntelliSense / hover tooltips, and because it can be
mechanically extracted into static HTML by tools such as DocFX or Sandcastle
without any further annotation work.

Every public type, public/internal member, and non-trivial private helper
carries at minimum a \texttt{<summary>} tag. Methods that take parameters or
return a meaningful value additionally carry \texttt{<param>} and
\texttt{<returns>} tags. Cross-references between related types use the
\texttt{<see cref="...">} tag so that IntelliSense renders clickable links.

\subsection*{Coverage}

\begin{table}[H]
\centering
\begin{tabular}{|l|r|}
\hline
\textbf{Metric} & \textbf{Count} \\ \hline
Custom C\# scripts in repository & 33 \\ \hline
Gameplay scripts (\texttt{Main Game/Scripts/}) & 22 \\ \hline
Menu / settings scripts (\texttt{Main Game/Menu/}) & 6 \\ \hline
Editor tooling (\texttt{Editor/}) & 1 \\ \hline
EditMode unit-test scripts (\texttt{Tests/EditMode/}) & 4 \\ \hline
Total \texttt{<summary>} doc tags & 209 \\ \hline
Total \texttt{<param>} / \texttt{<returns>} tags & 141 \\ \hline
\end{tabular}
\caption{XML documentation coverage across the custom codebase}
\end{table}

The XML documentation pass was committed in commit \texttt{b720a97}
(``Add XML documentation to all 29 custom C\# scripts''); the four test
scripts were documented inline as they were written. Unity-generated assets
(\texttt{TextMesh Pro}, \texttt{InputSystem\_Actions}, \texttt{OutlineEffect},
\texttt{UI Toolkit}, \texttt{Settings}) are excluded --- they are third-party
boilerplate, not authored team code.

\section{Documentation Conventions}

\subsection*{Class-level summaries}

Every class begins with a multi-line \texttt{<summary>} block describing its
single responsibility and any collaborators. For example, \texttt{PlayerKnowledge}
begins:

\begin{verbatim}
/// <summary>
/// Per-player inference about who holds which card. Populated by
/// GameManagerScript during suggestion resolution. Read by AiPlayer; in
/// a future PR, an "auto-fill notepad" toggle will project this into the
/// human's NotepadGrid.
///
/// Conceptually three tables of facts about the cards this player has
/// observed:
///   1. Definitely held by X       (saw X reveal it, or it's in own hand).
///   2. Definitely NOT held by X   (X passed on a suggestion containing
///                                  the card).
///   3. Held by exactly one of {X, cards}  (X refuted a suggestion but you
///                                          didn't see which card).
/// </summary>
public sealed class PlayerKnowledge { ... }
\end{verbatim}

\subsection*{Method-level documentation}

Public methods document their summary, parameters, and return value. The
following excerpt from \texttt{RoomRules.cs} is representative:

\begin{verbatim}
/// <summary>
/// Returns the room ID at the other end of the secret passage from
/// <paramref name="roomID"/>, or -1 if no passage exists.
/// </summary>
/// <param name="roomID">The room to travel from.</param>
/// <returns>Destination room ID, or -1.</returns>
public static int GetPassageDestination(int roomID) =>
    PassageMap.TryGetValue(roomID, out int d) ? d : -1;
\end{verbatim}

\subsection*{Enum and struct documentation}

Enum members and struct fields are documented individually so IntelliSense
explains each option at the call site. From \texttt{Player.cs}:

\begin{verbatim}
/// <summary>Actions available at the start of a turn (before or instead
/// of rolling).</summary>
public enum TurnAction
{
    /// <summary>Roll the dice and move.</summary>
    Roll,
    /// <summary>Use a secret passage to an adjacent room.</summary>
    Passage,
    /// <summary>Make a suggestion without rolling (only valid when already
    /// in a room).</summary>
    Suggest
}
\end{verbatim}

\subsection*{Inline comments}

Inline (\texttt{//}) comments are reserved for non-obvious algorithmic steps,
Unity-specific gotchas (e.g.\ Edit Mode vs.\ Play Mode coroutine constraints in
the test helpers), and rationale that would otherwise be lost. Self-explanatory
code is left uncommented to avoid noise.

\section{Documentation Tooling}

\begin{itemize}
    \item \textbf{IDE integration:} Visual Studio, VS Code (with the C\# Dev
    Kit), and JetBrains Rider all surface XML doc comments as hover tooltips
    and parameter hints. No separate build step is required --- every
    contributor sees the documentation while editing.
    \item \textbf{Static HTML generation:} the comments are compatible with
    DocFX (\texttt{docfx.exe}) and Sandcastle, which can produce a browsable
    HTML reference site directly from the source tree. This was not run for
    submission but the source is ready for it.
    \item \textbf{Compiler-checked references:} \texttt{<see cref="...">} tags
    are validated by the C\# compiler when the project is built with
    \texttt{/doc} enabled, catching stale references during refactors.
    \item \textbf{Project-level documentation:} architectural narrative,
    sprint records, and design diagrams live in the \texttt{report/} tree
    (Markdown and PlantUML sources, rendered into this report).
\end{itemize}

\section{Codebase Structure}

\subsection*{Repository layout}

\begin{table}[H]
\centering
\begin{tabular}{|p{6.2cm}|p{8cm}|}
\hline
\textbf{Path} & \textbf{Contents} \\ \hline
\texttt{README.md} & Setup guide: cloning, opening the project in Unity, Smart
Merge configuration, commit / PR workflow conventions \\ \hline
\texttt{clue/} & Unity project root (open this folder in Unity Hub) \\ \hline
\texttt{clue/Assets/Main Game/} & Game scene, prefabs, scripts \\ \hline
\texttt{clue/Assets/Main Game/Scripts/} & 22 gameplay scripts (board,
deck, player, AI, UI manager, etc.) \\ \hline
\texttt{clue/Assets/Main Game/Menu/} & 6 menu scripts (controller,
settings, stepper, modal, room plaque) \\ \hline
\texttt{clue/Assets/Editor/} & Editor-only tooling
(\texttt{ProjectFileGenerator.cs}) \\ \hline
\texttt{clue/Assets/Tests/EditMode/} & Sprint-aligned unit tests
(\texttt{Sprint1Tests.cs}\,--\,\texttt{Sprint4Tests.cs}) \\ \hline
\texttt{clue/Assets/StreamingAssets/} & \texttt{cards.json} --- runtime
data contract for the deck \\ \hline
\texttt{documentation/} & All design / planning / testing / meeting documentation
sources (Markdown + PlantUML) \\ \hline
\texttt{report/} & LaTeX source of this submission report \\ \hline
\end{tabular}
\caption{Repository layout and documentation locations}
\end{table}

\subsection*{Gameplay scripts at a glance}

Each script in \texttt{Main Game/Scripts/} owns a single responsibility,
documented in its class summary:

\begin{table}[H]
\centering
\resizebox{\textwidth}{!}{%
\begin{tabular}{|l|p{10cm}|}
\hline
\textbf{Script} & \textbf{Responsibility (from class \texttt{<summary>})} \\ \hline
\texttt{GameManagerScript.cs} & Central game controller: main turn loop, player
spawning, weapon placement, suggestion / accusation resolution \\ \hline
\texttt{BoardScript.cs} & 25$\times$24 tile grid; single source of truth for
walkability and occupancy \\ \hline
\texttt{DeckScript.cs} & Loads cards from \texttt{cards.json}, builds the
murder envelope, deals hands \\ \hline
\texttt{Player.cs} & Abstract base: turn state machine and decision hooks \\ \hline
\texttt{HumanPlayer.cs} & Player subclass: routes decisions through
\texttt{UiManager} \\ \hline
\texttt{AiPlayer.cs} & Player subclass: BFS-directed movement and
knowledge-filtered suggestions \\ \hline
\texttt{TokenMover.cs} & Grid movement, room entry / exit, world-space
positioning \\ \hline
\texttt{RoomRules.cs} & Static helper: secret-passage map and door-tile
detection \\ \hline
\texttt{DiceRollScript.cs} & Rigidbody dice; result decoded from face
orientation \\ \hline
\texttt{UiManager.cs} & Singleton entry point for every player-facing prompt \\ \hline
\texttt{NotepadUI.cs} & 147-cell detective-notebook toggle grid \\ \hline
\texttt{PlayerKnowledge.cs} & Per-player card inference (Held / NotHeld /
disjunctions) \\ \hline
\texttt{GameStateScript.cs} & Playing / MainMenu / GameOver state machine \\ \hline
\texttt{SceneController.cs} & Persists player and AI count between menu and
game scenes \\ \hline
\texttt{CharacterData.cs} & ScriptableObject: per-character static metadata \\ \hline
\texttt{GameOverScreen.cs} & End-of-game restart / return-to-menu UI \\ \hline
\texttt{CameraFollow.cs} & Smooth camera follow for the active player \\ \hline
\texttt{BoardThemer.cs} / \texttt{Theme Changer.cs} & Runtime board / menu
theme switching \\ \hline
\texttt{BakeData.cs} / \texttt{BakedNavData.cs} & Pre-computed navigation data
for AI BFS \\ \hline
\texttt{DiceDebugLogger.cs} & Editor-only diagnostic logger for dice physics \\ \hline
\end{tabular}%
}
\caption{Gameplay scripts and their documented responsibilities}
\end{table}

\section{Setup and Usage Guide}

The repository \texttt{README.md} is the canonical setup document. The key
steps it covers are:

\begin{enumerate}
    \item Clone the repository to a known location.
    \item Install the Unity version pinned in the project (Unity Hub will
    prompt if a different version is opened).
    \item In Unity Hub, click \textit{Add} $\to$ select the \texttt{clue}
    folder.
    \item Open \textit{Assets / Main Game / MenuScene} 
    press \textit{Play}.
\end{enumerate}

\noindent The README additionally documents the team's commit-message
expectations (testing notes, design rationale, collaboration evidence, risk
assessment, task attribution) and the branching strategy (feature branches per
sprint, merged via reviewed pull requests).

\newpage
